% Draft (20 Oct 2026) of the new coupling text for Chapter 4, replacing the scaled scheme of
% Section \ref{sec:ch5-composition}. Not included in the build yet. Plan: thesis/CH4_REWRITE_PLAN.md.
% Decisions: MEETING_2026-10-20.md. Numbers marked [..] are filled in from the runs.

\section{Coupling the point model to the field}
\label{sec:ch5-composition}

\paragraph{The point model as published.}
Klempt et al.~\cite{Klempt2026ContinuumBacterialGrowth} formulate their model on a
material point, with $\phi_0$ as an auxiliary variable that enforces $\sum_i\phi_i+\phi_0=1$
and is not to be read as a physical void. I therefore do not change the point model: its
equations~\eqref{eq:ch5-point}, its parameters and its constraint stay as published, and it is
called at every Gauss point with given inputs and returns its state. The inputs are the nutrient $c^*$,
the antibiotic $\alpha^*$ and the start values; the outputs are $\phi_i$ and $\psi_i$. The
field keeps its own amount $\phi_{3D}$, and the point model keeps its own state, which is carried
from one increment to the next as state variables of the Gauss point.

\paragraph{From the field to the point model.}
The element solves the nutrient $c(\mathbf x,t)$ of Klempt et al.\ 2024, held at given sources
and diffusing through the domain. The point model needs $c^*$ at its own, much smaller time
steps. Between the increments $t_n$ and $t_{n+1}$ I interpolate linearly,
\begin{equation}
c^*(t)=c^*_0\Bigl[\tfrac{t_{n+1}-t}{\Delta t}\,c_n+\tfrac{t-t_n}{\Delta t}\,c_{n+1}\Bigr],
\qquad t_n\le t\le t_{n+1},
\label{eq:ch5-cstar}
\end{equation}
with $c_n,c_{n+1}$ the nutrient at the Gauss point and $c^*_0$ the nutrient level of the point
model's own examples. Klempt et al.\ 2026 already use a time-dependent $c^*$ (their four-species
case~3), so the point model accepts such an input without change. The model of Klempt et al.\
2024 has no antibiotic. The element solves a second nutrient for its second species, which I do
not use; I use this field for the antibiotic $a(\mathbf x,t)$, with its own sources and the same
diffusion, and pass $\alpha^*(t)=\alpha^*_0\,a$ to the point model, interpolated as
in~\eqref{eq:ch5-cstar}.

\paragraph{From the point model to the field.}
Only living biomass consumes nutrient. In Eq.~35 of Klempt et al.\ 2024 the consumption is
written $g\phi$ with $\phi=\phi_{3D}$; the implementation uses $g\phi_{3D}c$, and the printed form is a misprint.
Throughout, $\phi_{3D}$ denotes the field of Klempt et al.\ 2024 and $\phi_i$ the point model. I scale
the consumption by the living biomass at the Gauss point,
\begin{equation}
0=d\,\nabla^2c-g\,\phi_{3D}\Bigl(\sum_{i=1}^N\bar\phi_i\Bigr)c,\qquad
\bar\phi_i=\phi_i\psi_i,
\label{eq:ch5-consumption}
\end{equation}
with $\bar\phi_i$ the living volume fraction of species $i$ in the notation of Klempt et al.\ 2026.
A dead biofilm consumes nothing; the void $\phi_0$ has no living fraction and is not part of the sum.
The point model normalises its volume, $\sum_i\phi_i+\phi_0=1$, so $\bar\phi_i\in[0,1]$ is
the living volume fraction of species $i$ at the point, and $\phi_{3D}\sum_i\bar\phi_i\le\phi_{3D}$ is
the living volume fraction in the field. The antibiotic is
taken up in the same way, $0=d_a\nabla^2a-g_a\phi_{3D}\bigl(\sum_i\bar\phi_i\bigr)a$. The growth law~\eqref{eq:ch5-eq36} is
unchanged: the inflation $\alpha$ measures how much biomass has been produced, whether alive or
dead. The spreading of $\phi_{3D}$ follows Eq.~34 with its front term, which depends on $c$; this
is the path by which the point model reaches the growth and the stress. $\sum_i\bar\phi_i$ is taken from the
previous increment.

\paragraph{Time.}
Klempt et al.\ 2024 normalise time with a reference time chosen for the problem at hand. I choose
$T^*=1$ as the duration covered by the point model's data, so that both models run on the same
time: the point model advances $s\,\Delta T^*$ per increment with $s$ the length of that window in
its own units, not a fitted value. For the synthetic data of Fritsch et al.\
\cite{Fritsch2025BayesianMicrofilms}, case~I (1000 steps of $10^{-4}$), $s=0.1$. With this choice
the coupled run reproduces the point model alone: [replay numbers].

\paragraph{Start values.}
The point model starts where the field first reaches $\phi_{\min}=0.01$. Where a seed is placed,
the start values are those of the seed (one species, or a mixture). A Gauss point that the
spreading biofilm reaches later starts from the $\phi_{3D}$-weighted mean state of the points it
grows from. This is an assumption: the point model has no transport between points, so species
meet only through this rule. [sensitivity to $\phi_{\min}$].

\paragraph{Assumptions.} The coupling rests on four choices that neither paper fixes, each
with its sensitivity in Section~\ref{sec:ch5-results}: the linear interpolation of $c^*$ and
$\alpha^*$ in time, $L$ from the previous increment, the start values of newly reached points,
and $\phi_{\min}$. The time link is no longer one of them.
